\section{Introduction}

The hippocampus is an important subcortical area, commonly viewed as a system capable of both rapid encoding single episodes in joint representations and retrival from partial cues.
Another perspective considers its neural dynamics to be complementary to the slower cortical learning \cite{teyler1986, mcclelland1995}.
A performant memory system must thus be able to differentiate between similar observations, pre-process and recall the stored states, and present them in a format usable by downstream circuits.
The first feature is pivotal for pattern separations, while the second one concerns more pattern completion and attractor-like recall.
Here, the emphasis is placed on the third one, focusing on the representational coordinates in which the retrieve content is projected.
This motivates an interpretation of the hippocampal sub-field CA1 where it integrates successive experiences within a developed latent neural space.

The neuroanatomy of the hippocampal formation is organized in subfield covering a distribution of computational functions.
The main afferent region is the entorhinal cortex (EC), whose projections are generally distincted into the medial EC (MEC), often associated to grid cells and spatial metrics from path integration, and the lateral EC (LEC), which convey more predominately local sensory, object, and event-releated
information.
Nonetheless, both regions express a mixture of spatial and nonspatial features, with their corresponding contributions being influenced by behavioural and task context \cite{hargreaves2005, henriksen2010, deshmukh2011}.
The principal trisynaptic loop delivers inputs from the entorhinal layer II through the dentate gyrus (DG), then CA3 and CA1; whereas entorhinal layer III constitures a direc temporoammonic pathway to CA1.
In the processes of context-addressable retrieval, auto-associative models emphasize recurrent CA3 circuitry \cite{marr1971}, possibly with a spatially distributed organization \cite{li2024, chandra2025}.
Aligned with this perspective, the selective silencing of CA3 NMDA-receptor function affects pattern completion from partial cues, while CA3 and CA1 ensembles respond differently to environmental changes \cite{nakazawa2002,leutgeb2004}.
These approaches constitute mechanisms by which learned representation could be recalled, but they do not by themsleves specify the internal neural coordinates in which the expressed activity should be manifested in the hippocampal response.

Representation in CA1 population can emerge and reorganiza rapidly.
Behavioural-timescale synaptic plasticity (BTSP) is a modeled learning rule observed to modify the place fields of CA1 pyramidal neurons by integrating over the interaction of instructive events and CA3 projections \cite{bittner2017, milstein2021}.
The activity of entorhinal layer III is typically implicated as the instructive or teaching-related signal to incorporate within CA1 populations \cite{grienberger2022}.
Experimental observations and modeling further link BTSP dynamics to trial-by-trial shifts of place fiels and the stabilization of task-oriented spatial encodings \cite{madar2025, vaidya2025}..
The relevance of BTSP thus lies both in the involvement in place cells formation as well as rapid adaptation of spatial representation upon contigencies changes.
Computational work additionally show how BTSP-like plasticity rules can manifest content-addressable storage dynamics \cite{wu2025}.
The present study abstracts this pattern as an entorhinal IS used to update CA3-CA1 synapic weights; it does not propose an account of eligibility traces, dendritic plateau potentials, or other biophysical details of BTSP.

Another property exhibited by CA1 activity is sensitivity to context.
Shifts in environmental or cue configurations can affect local firing rates, field shapes, and the morphology of ensembles, producing phenomena such as rate and global remapping \cite{muller1987, leutgeb2005}.
The direct entorhinal projections deliver spatial and nonspatial information that can be jointly determine these changes.
In the model proposed in this work, distinct MEC and LEC-like input partitions offer a controlled way to represent those channels, aimed at testing pathwy-selective perturbation.

Our main proposal adds a representational conditions to this circuit motif.
If rapid plasticity at CA1 synapses encode patterns whicle the downdestream readout stays fixed, the pattern may retain information regarding the stimulus but being functionally useless if it falls outside the neural activity subspace or internal coordinates to which the readout decoder is tuned.
Related constraints have been osberved outside the hippocampus.
For instance, in brain-computer-interface experiments, motor-cortical perturbations that stayed within previous neural manifoldwere learned mroe efficiently than those necessitating activity information outside that manifold \cite{sadtler2014}.
Addtionally, short-term adaptation could also take place in the form of re-activation of former cortical activity together with new behavioural evidence \cite{golub2018}.
Homeostatic and stability dynamics at the population level can help preserving behavioural representation despite consistent changes in the subsets of involved neurons \cite{gallego2020}.
Together, these findings suggests how established downstream mappings may constrain the recruitment and consolidation of new learned populations states.
Here, we adopt the narrower term \emph{readout compatibility} to specify the property of the newly instructed CA1 state of remaining within the activity subspace that a learned output mapping can decode correctly.
Readout compatibility is therefore considered a functional criteria, meanwhile manifold preservation is a possible process by which such correspondance could be maintained.

In order to imprint prior latent structure in the direct EC connections, an EC-CA1-EC autoencoder is pre-trained before the application of plasticity to store new associations. 
Particularly, this is a computational trick relying on backpropagation optimization applied to a separate simulation phase, that does not claim neurobiological resemblance.
At an interpretation level, this pre-trained mapping is meant to encapsulate the accumulation of information through development, lifelong experiences, and relatively stable connections among entorhinal and hippocampal subfields \cite{mcclelland1995, chandra2025, wen2024}. 
In deep learning, this is a common technique used for extracting principal input features early on, before the proper learning loop is started with a possibly different simulation setup \cite{yosinski2014}.
Freezing the decoder, namely the CA1-EC weights, after pre-training turns the reconstruction goal into a task in which new CA1 assemblies promped by joint CA3 and EC inputs must remain decipherable by the readout, effectively leveraging and going beyond the learned encoder.
Such paradigm could in theory be able to support stable decoding over ongoing experiences, although long-term memorization is not within the scope of this work.

The model used here is deliberately minimal.
An entorhinal layer EC follows two routes: the autoencoding loop with CA1-EC being the latent layer, and a longer path through CA3-CA1-EC.
CA3 activity is effectively used as a retrival key used to modify the CA3-CA1 synapses with a BTSP-inspired learning paradigm, such to match the decodability performance achieved by the autoencoder pretraining.
Two plasticity variants are used, an instructive-driven, based on a more direct synaptic overwrite, and a bounded error-rule, which instead used the difference between instructive and recalled CA1 activity.
The simulations made intend to assess whether stimuli reconstruction depends on the compatibility of internal neural coordinates, whether shifts in cue-sequences induce reproducible CA1 states, what are the differential computational contributions of the two updates rules, and how they differ under
matched conditions.
Futher, it is tested how recall is affected by degraded input observation, and whether the number of distinct cue-pairs affect retrieval accuracy.
Finally, the goal of the model is to propose structural and operational representational conditions and exposes their consequences in simplified settings, without claiming to follow strict biological constraints. 




